(a) In a noetherian space, the finite unions of members of an open cover of an open subset have a maximal member; it must be the whole subset. Conversely, the union of an ascending chain of open subsets is quasi-compact, so is covered by finitely many members of the chain and equals one of them. Hence the chain stabilizes.

(b) Refine an open cover of $\operatorname{Spec}A$ by distinguished opens $D(f_i)$. Since no prime contains all the $f_i$, they generate the unit ideal. An expression $1=\sum_{i=1}^na_if_i$ gives a finite subcover. For $A=k[x_1,x_2,\ldots]$, however, the open subset $\bigcup_iD(x_i)$ has no finite subcover by these opens: $(x_j:j\ne i)$ belongs only to $D(x_i)$. Thus the spectrum need not be noetherian.

(c) Every open subset is a union $\bigcup_iD(f_i)$. The ideal generated by the $f_i$ is generated by finitely many of them, so this union is finite and hence quasi-compact. Apply (a).

(d) The ring $k[x_1,x_2,\ldots]/(x_1^2,x_2^2,\ldots)$ is not noetherian, since $(x_1)\subsetneq(x_1,x_2)\subsetneq\cdots$. Its nilradical is $(x_1,x_2,\ldots)$ and its reduced ring is $k$, so its spectrum has one point.
